\section{Related Work}

\paragraph{Federated learning under statistical and system heterogeneity.}
Federated learning (FL) studies how multiple clients collaboratively train a model without sharing raw data, with FedAvg establishing the canonical protocol of repeated local optimization followed by server aggregation \cite{mcmahan2016fedavg}. Early system work further clarified the practical communication and deployment constraints of cross-device FL \cite{konevcny2016communication,bonawitz2016practical}. A central difficulty is statistical heterogeneity: when client data are non-IID, local training induces client drift, which can substantially slow or even destabilize optimization. A major line of work therefore designs heterogeneity-aware objectives or corrections. FedProx stabilizes local training by adding a proximal term \cite{li2018fedprox}, while SCAFFOLD uses control variates to correct drift and achieves substantially improved convergence guarantees under heterogeneous data \cite{karimireddy2019scaffold}. Related analyses further showed that naive aggregation can suffer from objective inconsistency when clients perform unequal amounts of local work \cite{wang2020objective}, motivating more principled algorithm design for heterogeneous participation. Beyond first-order averaging, Newton-type and primal--dual formulations such as FedDANE and more recent localized primal--dual methods extend FL to broader objective classes and heterogeneous regimes \cite{li2019feddane,tyou2024localgecl,zhou2025fedcanon}. Personalization methods, including meta-learning, Moreau-envelope formulations, representation sharing, and distillation-based approaches, improve client-specific performance but typically introduce extra local state, auxiliary models, or additional communication rounds \cite{fallah2020perfedavg,dinh2020pfedme,collins2021exploiting,jeong2018federated}. Overall, existing FL methods have made substantial progress on non-IID robustness and participation heterogeneity, but they predominantly optimize statistical efficiency and communication cadence, leaving client memory footprint largely untreated.

\paragraph{Communication-efficient optimization and compressed federated/distributed training.}
Communication is one of the dominant bottlenecks in FL and distributed learning, which has motivated a large literature on quantization, sparsification, local updates, and decentralized training. Classical quantization methods such as QSGD reduce transmitted bits while retaining unbiasedness \cite{alistarh2016qsgd}. Orthogonal to compression, local-update methods reduce synchronization frequency; representative analyses include Local SGD and the broader Cooperative SGD framework \cite{lin2018localsgd,wang2018cooperative}. In practice, however, aggressive biased compressors often require explicit error compensation to avoid optimization bias. Modern error-feedback methods, especially EF21, provide a cleaner and stronger theoretical foundation for compressed optimization, while subsequent work refined the understanding of sparsification as a total-error trade-off rather than a pure bandwidth problem \cite{richtarik2021ef21,sahu2021rethinking}. Momentum and acceleration further improve communication efficiency, and compressed methods approaching lower-bound-optimal behavior have been developed in both distributed and federated settings \cite{li2020accelerationcompressed,huang2022lowerbounds}. Recent FL-specific work combines drift correction and compression more tightly, for example by extending SCAFFOLD-style control variates to compressed communication \cite{huang2023stochastic}. Related algorithms such as DASHA and BEER achieve strong oracle or communication complexity guarantees under compression \cite{tyurin2022dasha,zhao2022beer}, while low-rank communication schemes such as PowerSGD provide practical throughput benefits in large-scale training \cite{vogels2019powersgd}. A complementary direction removes the central server and studies decentralized protocols, including D-PSGD, asynchronous decentralized SGD, and D$^2$, which mitigate synchronization bottlenecks and can reduce or even remove the adverse effect of data heterogeneity in certain regimes \cite{lian2017dpsgd,lian2017async,tang2018d2}. More recent analyses of gradient tracking have sharpened convergence guarantees for decentralized learning and enabled robust local-update variants over dynamic networks \cite{koloskova2022improved,ghiasvand2024robust}. Despite this rich line of work, the optimization target remains communication volume; these methods usually still assume the ability to compute and store dense gradients, model states, and optimizer statistics on each client.

\paragraph{Memory-efficient training and the missing joint perspective.}
A separate literature targets training-time memory usage. In large-model optimization, GaLore showed that projecting gradients into a low-rank subspace can drastically reduce optimizer-state memory while retaining strong empirical performance \cite{zhao2024galore}. Yet recent theory has also revealed a key caveat: naive stochastic subspace optimization can fail, motivating provably convergent subspace variants and a more careful treatment of noise, projection bias, and optimization geometry \cite{he2024subspace}. Other memory-oriented methods reduce state or activation cost through optimizer redesign and projection mechanisms \cite{horvath2024frugal}. In the FL setting, memory efficiency has only recently begun to attract attention, with approaches based on split execution for heterogeneous mobile devices, progressive layer freezing, sparse updates, or low-memory device adaptation \cite{chen2025splitfl,wu2024progressive,tinyfel2025}. While these methods demonstrate that memory is a first-class deployment bottleneck, most are designed primarily as systems heuristics and provide limited optimization theory under heterogeneous local data.

The gap becomes sharper when comparing the memory and communication literatures side by side. Communication-efficient FL methods typically preserve standard Euclidean local optimization and only compress transmitted messages \cite{alistarh2016qsgd,richtarik2021ef21,huang2023stochastic}; memory-efficient methods instead modify the local optimization geometry itself through projection, freezing, sparsity, or split execution \cite{zhao2024galore,he2024subspace,wu2024progressive,chen2025splitfl}. Once client updates are produced by non-Euclidean or subspace-constrained local steps, however, the assumptions underpinning classical drift-correction analyses no longer apply directly. Existing FL theory therefore does not yet characterize how to jointly control \emph{client drift}, \emph{compression error}, and \emph{memory-induced projection bias} in one framework. This missing joint treatment is particularly consequential in heterogeneous FL, where practical clients are simultaneously limited by communication bandwidth, optimizer-state memory, activation memory, and irregular participation.

In summary, prior work has separately advanced heterogeneity-robust FL \cite{li2018fedprox,karimireddy2019scaffold,wang2020objective}, communication-efficient optimization \cite{alistarh2016qsgd,richtarik2021ef21,li2020accelerationcompressed,huang2022lowerbounds,tyurin2022dasha,zhao2022beer}, decentralized and gradient-tracking methods \cite{lian2017dpsgd,tang2018d2,koloskova2022improved,ghiasvand2024robust}, and memory-efficient training \cite{zhao2024galore,he2024subspace,horvath2024frugal,wu2024progressive,chen2025splitfl,tinyfel2025}. What remains largely open is a principled FL algorithm that simultaneously achieves \emph{memory efficiency} and \emph{communication efficiency}, while preserving rigorous convergence guarantees under non-IID data and partial participation. Our work is motivated precisely by this missing intersection.
